% GENERATED FILE. Edit canonical lessons, then run npm run generate:books.
% canonical-source-hash: fnv1a64:a5fc7549e18716ca
% canonical-lessons: FR-C484-ambulance, FR-C484-secouriste, FR-C484-brancard, FR-C484-piqure, FR-C484-salle-d-attente

\chapter{Aux urgences}
\label{ch:fr-aux-urgences}
\hlchaptermodality{484}

\begin{chapteropening}
\textbf{By the end of this chapter:} \emph{I can talk about emergency care.}

Complete the last lesson of chapter 484: I can talk about emergency care.
\end{chapteropening}

\section[l'ambulance]{l'ambulance — ambulance}

\label{lesson:FR-C484-ambulance}



\begin{quote}
Before the new one: say the French for dizziness, then the French for to digest.
\end{quote}

\subsection*{You'll want to know: l'ambulance}
\textbf{l'ambulance} — \textquotedblleft{}ambulance\textquotedblright{}.

\begin{quote}
L'ambulance est arrivée en dix minutes.
\end{quote}

\begin{quote}
\emph{The ambulance arrived in ten minutes.}
\end{quote}

Feminine: \emph{une ambulance}. \emph{Appeler une ambulance} — to call an ambulance.

\begin{cousinweb}
\emph{l'ambulance} is built on \emph{ambulant} \textquotedblleft{}moving about\textquotedblright{} (EtymDB), from Latin \emph{ambulāre} \textquotedblleft{}to walk\textquotedblright{} (Lewis \& Short): first a field hospital that moved with an army.
\end{cousinweb}

\begin{grammarlens}[title={What you've built}]
One more word for telling, explaining and describing at B1.
\end{grammarlens}

\subsection*{Your turn}
\begin{itemize}
\raggedright
  \item \emph{Say it:} \emph{l'ambulance}
  \item \emph{Say it:} \emph{l'ambulance}, once more
  \item \emph{Say it:} a sentence of your own with \emph{l'ambulance}
  \item \emph{Recall:} say the French for dizziness, then the French for to digest, then say \emph{l'ambulance} again
\end{itemize}

\begin{tcolorbox}[breakable,colback=teal!4,colframe=teal!35!black,title={Before you move on}]
What does l'ambulance mean? (\textquotedblleft{}Ambulance\textquotedblright{}.) Say it once more.
\end{tcolorbox}
\section[le secouriste]{le secouriste — first-aider, rescue worker}

\label{lesson:FR-C484-secouriste}



\begin{quote}
Before the new one: say the French for to digest, then the French for ambulance.
\end{quote}

\subsection*{You'll want to know: le secouriste}
\textbf{le secouriste} — \textquotedblleft{}first-aider, rescue worker\textquotedblright{}.

\begin{quote}
Un secouriste s'est occupé du blessé.
\end{quote}

\begin{quote}
\emph{A first-aider took care of the injured man.}
\end{quote}

A woman is \emph{la secouriste}. \emph{Les premiers secours} — first aid.

\begin{cousinweb}
\emph{le secouriste} is built on \emph{secours} \textquotedblleft{}help\textquotedblright{}, which goes back to Latin \emph{succursus}, from \emph{succurrere} \textquotedblleft{}to run to help\textquotedblright{}, built on \emph{currere} \textquotedblleft{}to run\textquotedblright{} (Lewis \& Short).
\end{cousinweb}

\begin{grammarlens}[title={What you've built}]
One more word for telling, explaining and describing at B1.
\end{grammarlens}

\subsection*{Your turn}
\begin{itemize}
\raggedright
  \item \emph{Say it:} \emph{le secouriste}
  \item \emph{Say it:} \emph{le secouriste}, once more
  \item \emph{Say it:} a sentence of your own with \emph{le secouriste}
  \item \emph{Recall:} say the French for to digest, then the French for ambulance, then say \emph{le secouriste} again
\end{itemize}

\begin{tcolorbox}[breakable,colback=teal!4,colframe=teal!35!black,title={Before you move on}]
What does le secouriste mean? (\textquotedblleft{}First-aider, rescue worker\textquotedblright{}.) Say it once more.
\end{tcolorbox}
\section[le brancard]{le brancard — stretcher}

\label{lesson:FR-C484-brancard}



\begin{quote}
Before the new one: say the French for ambulance, then the French for first-aider.
\end{quote}

\subsection*{You'll want to know: le brancard}
\textbf{le brancard} — \textquotedblleft{}stretcher\textquotedblright{}.

\begin{quote}
Ils ont porté le joueur sur un brancard.
\end{quote}

\begin{quote}
\emph{They carried the player on a stretcher.}
\end{quote}

Also called \emph{une civière}.

\begin{cousinweb}
\emph{le brancard} is built on a northern dialect form of \emph{branche} \textquotedblleft{}branch\textquotedblright{} with \emph{-ard}: first two poles for carrying. \emph{Branche} comes from Late Latin \emph{branca} \textquotedblleft{}paw\textquotedblright{}.
\end{cousinweb}

\begin{grammarlens}[title={What you've built}]
One more word for telling, explaining and describing at B1.
\end{grammarlens}

\subsection*{Your turn}
\begin{itemize}
\raggedright
  \item \emph{Say it:} \emph{le brancard}
  \item \emph{Say it:} \emph{le brancard}, once more
  \item \emph{Say it:} a sentence of your own with \emph{le brancard}
  \item \emph{Recall:} say the French for ambulance, then the French for first-aider, then say \emph{le brancard} again
\end{itemize}

\begin{tcolorbox}[breakable,colback=teal!4,colframe=teal!35!black,title={Before you move on}]
What does le brancard mean? (\textquotedblleft{}Stretcher\textquotedblright{}.) Say it once more.
\end{tcolorbox}
\section[la piqûre]{la piqûre — injection; sting, bite}

\label{lesson:FR-C484-piqure}



\begin{quote}
Before the new one: say the French for first-aider, then the French for stretcher.
\end{quote}

\subsection*{You'll want to know: la piqûre}
\textbf{la piqûre} — \textquotedblleft{}injection; sting, bite\textquotedblright{}.

\begin{quote}
L'infirmière m'a fait une piqûre au bras.
\end{quote}

\begin{quote}
\emph{The nurse gave me an injection in the arm.}
\end{quote}

Also an insect bite: \emph{une piqûre de moustique}. Also spelt \emph{piqure}.

\begin{cousinweb}
\emph{la piqûre} is built on \emph{piquer} \textquotedblleft{}to prick\textquotedblright{} with \emph{-ure}. \emph{Piquer} goes back to a Vulgar Latin \emph{piccāre}, probably an expressive word.
\end{cousinweb}

\begin{grammarlens}[title={What you've built}]
One more word for telling, explaining and describing at B1.
\end{grammarlens}

\subsection*{Your turn}
\begin{itemize}
\raggedright
  \item \emph{Say it:} \emph{la piqûre}
  \item \emph{Say it:} \emph{la piqûre}, once more
  \item \emph{Say it:} a sentence of your own with \emph{la piqûre}
  \item \emph{Recall:} say the French for first-aider, then the French for stretcher, then say \emph{la piqûre} again
\end{itemize}

\begin{tcolorbox}[breakable,colback=teal!4,colframe=teal!35!black,title={Before you move on}]
What does la piqûre mean? (\textquotedblleft{}Injection; sting, bite\textquotedblright{}.) Say it once more.
\end{tcolorbox}
\section[la salle d'attente]{la salle d'attente — waiting room}

\label{lesson:FR-C484-salle-d-attente}



\begin{quote}
Before the new one: say the French for stretcher, then the French for injection.
\end{quote}

\subsection*{You'll want to know: la salle d'attente}
\textbf{la salle d'attente} — \textquotedblleft{}waiting room\textquotedblright{}.

\begin{quote}
Nous avons passé deux heures dans la salle d'attente.
\end{quote}

\begin{quote}
\emph{We spent two hours in the waiting room.}
\end{quote}

At a doctor's, a station or a hospital.

\begin{cousinweb}
\emph{la salle d'attente} joins \emph{salle} \textquotedblleft{}room\textquotedblright{}, from Frankish \emph{sal} \textquotedblleft{}a hall\textquotedblright{}, and \emph{attente} \textquotedblleft{}waiting\textquotedblright{}, built on \emph{attendre}, inherited from Latin \emph{attendere} \textquotedblleft{}to stretch towards\textquotedblright{}, from \emph{tendere} (Lewis \& Short).
\end{cousinweb}

\begin{grammarlens}[title={What you've built}]
One more word for telling, explaining and describing at B1.
\end{grammarlens}

\subsection*{Your turn}
\begin{itemize}
\raggedright
  \item \emph{Say it:} \emph{la salle d'attente}
  \item \emph{Say it:} \emph{la salle d'attente}, once more
  \item \emph{Say it:} a sentence of your own with \emph{la salle d'attente}
  \item \emph{Recall:} say the French for stretcher, then the French for injection, then say \emph{la salle d'attente} again
\end{itemize}

\begin{tcolorbox}[breakable,colback=teal!4,colframe=teal!35!black,title={Before you move on}]
What does la salle d'attente mean? (\textquotedblleft{}Waiting room\textquotedblright{}.) Say it once more.
\end{tcolorbox}
